IJM is~configured at~three layers: a~pair of~JSON files that describe
the~cluster, an~HTTP submission API for~jobs, and~a~handful of~runtime
environment variables.  This section documents each in~turn.

\section{Cluster configuration}

The cluster topology is~loaded at~startup from two JSON files, whose
paths are configurable via~environment variables.

\subsection{Nodes (\texttt{nodes\_config.json})}

Each entry describes a~compute node and~its GPU resources.  The~example
below is~the~deployment used throughout
Chapter~\ref{ch:e2e}~(\texttt{config/nodes\_config.tunnel.json}); the
file checked in~at~\texttt{config/nodes\_config.json} is~a~mock cluster
used by~the~unit tests.

\begin{lstlisting}[language={}, numbers=none]
[
  {
    "id": "matemagician",
    "isForProfiling": true,
    "workerUrl": "http://localhost:8001",
    "cost": 0.10,
    "resources": [
      { "gpu_type": "QuadroP600", "gpu_count": 2 }
    ]
  },
  {
    "id": "polimi-gpu",
    "isForProfiling": true,
    "workerUrl": "http://localhost:8002",
    "cost": 0.30,
    "resources": [
      { "gpu_type": "A40", "gpu_count": 2 }
    ]
  }
]
\end{lstlisting}

\begin{table}[H]
\centering
\begin{tabularx}{\textwidth}{@{}l l X@{}}
\toprule
\textbf{Field} & \textbf{Type} & \textbf{Description} \\
\midrule
\texttt{id}                     & string         & Unique node identifier \\
\texttt{isForProfiling}         & bool           & If \texttt{true}, the~node can host profiling runs.  Production jobs may still land on~it; profiling jobs may only land on~it \\
\texttt{workerUrl}              & string         & HTTP base URL of~the~node's worker, reached through an~SSH tunnel.  The~API dispatches by~posting \texttt{/jobs/\{id\}/run} and~\texttt{/jobs/\{id\}/stop} here; every node must be~backed by~a~worker \\
\texttt{cost}                   & float          & Idle cost of~the~node ($\$$/h), used by~the~optimiser \\
\texttt{resources}              & array          & GPU resources attached to~the~node \\
\texttt{resources[].gpu\_type}  & string         & GPU model name (e.g.~\texttt{A40}, \texttt{L40S}, \texttt{QuadroP600}) \\
\texttt{resources[].gpu\_count} & int            & Number of~GPUs of~this type \\
\bottomrule
\end{tabularx}
\caption{Node configuration fields.}\label{tab:node-config}
\end{table}

GPU type and~count are declared explicitly rather than auto-detected.
This admits mixed-GPU nodes that ANDREAS did not support (e.g.~a~node
with~both A40 and~L40S GPUs), and~it~lets the~operator declare nodes
that are not physically attached to~the~API host~--- the~API only ever
talks to~workers over HTTP.

\subsection{Energy costs (\texttt{gpu\_energy\_costs.json})}

A~lookup table that maps a~GPU type and~bundle size to~an~hourly energy
cost (\$/h).  The~scheduler consults this table when scoring candidate
placements.

\begin{lstlisting}[language={}, numbers=none]
{
  "A40":        { "1": 0.30, "2": 0.55, "4": 1.05 },
  "L40S":       { "1": 0.40, "2": 0.75, "4": 1.40 },
  "QuadroP600": { "1": 0.03, "2": 0.055 },
  "Blackwell":  { "1": 1.00, "2": 1.85, "4": 3.50 }
}
\end{lstlisting}

\section{Job submission}

\textbf{Endpoint:} \texttt{POST /jobs}

\subsection{Request body}

\begin{lstlisting}[language={}, numbers=none]
{
  "job_id": "LSTM-small",
  "dockerImage": "ijm-lstm-small:dev",
  "scriptPath": "train.py",
  "directoryToMount": "/data/my-project",
  "Priority": 3,
  "deadline": "2026-03-20T18:00:00",
  "batchSize": 64,
  "profilingEpochsNo": 3,
  "epochsTotal": 20
}
\end{lstlisting}

\begin{table}[H]
\centering
\begin{tabularx}{\textwidth}{@{}l l l X@{}}
\toprule
\textbf{Field} & \textbf{Type} & \textbf{Required} & \textbf{Description} \\
\midrule
\texttt{job\_id}            & string       & yes & Reusable job type identifier; profile measurements are correlated across submissions sharing the~same \texttt{job\_id} \\
\texttt{dockerImage}        & string       & yes & Docker image to~execute \\
\texttt{scriptPath}         & string       & no  & Path to~the~training script (overrides Docker CMD with~\texttt{python3 <scriptPath>}) \\
\texttt{directoryToMount}   & string       & no  & Host directory persisted on~the~job row (ANDREAS-compatibility field); not~mounted by~the~current worker, which mounts only checkpoints, the~runs directory, and~the~shared dataset cache \\
\texttt{command}            & string[]     & no  & Command and~arguments (defaults to~the~image entrypoint; overridden by~\texttt{scriptPath}) \\
\texttt{Priority}           & int 1--5     & no  & Job priority (default~3) \\
\texttt{deadline}           & datetime     & no  & Job deadline (ISO 8601) \\
\texttt{batchSize}          & int          & no  & Batch size passed as~the~\texttt{BATCH\_SIZE} env var \\
\texttt{profilingEpochsNo}  & int $\geq 1$ & no  & Epochs per~profiling run (default~3) \\
\texttt{epochsTotal}        & int $\geq 1$ & no  & Total training epochs (default~20) \\
\bottomrule
\end{tabularx}
\caption{Job creation request fields.}\label{tab:job-create}
\end{table}

Each submission receives a~unique UUID instance ID (\texttt{id}) for
lifecycle tracking, separate from the~\texttt{job\_id} type identifier.
When \texttt{scriptPath} is~provided the~container runs
\texttt{python3 <scriptPath>}; otherwise the~image's default CMD is~used.
Unlike ANDREAS, \texttt{Priority} and~\texttt{deadline} are optional,
and~an~additional \texttt{command} field gives flexibility for
non-standard workloads.

\subsection{Response (201 Created)}

The response carries the~full job object, including the~scheduler's
initial decision: assigned node, GPU configuration, and~whether the
first run will be~a~profile or~a~standard execution.

\begin{lstlisting}[language={}, numbers=none]
{
  "id": "a1b2c3d4-...",
  "job_id": "LSTM-small",
  "image": "ijm-lstm-small:dev",
  "command": ["python3", "train.py"],
  "status": "QUEUED",
  "is_profiling_run": true,
  "assigned_node": "polimi-gpu",
  "assigned_gpu_config": { "A40": 1 },
  "priority": 3,
  "epochs_total": 20,
  "profiling_epochs_no": 3,
  "created_at": "2026-03-18T10:00:00Z",
  "updated_at": "2026-03-18T10:00:00Z",
  ...
}
\end{lstlisting}

\subsection{Job image requirements (the trainer contract)}
\label{sec:job-contract}

The fields above describe what the~API \emph{stores}; this subsection is~the
\emph{specification} a~container author implements so~a~job is~not~merely
runnable but~\emph{stoppable, resumable, and~migratable}.  A~job image is~a
container whose entrypoint instantiates a~subclass of~\texttt{BaseTrainer}
(\texttt{runtime/base.py}) and~calls \texttt{.train()}; subclassing \emph{is}
the~contract, and~in~return the~stop/resume/migrate machinery
of~\S\ref{sec:checkpoints} comes for~free.

\paragraph{What you implement.}  A~subclass provides up~to~three hooks
(Table~\ref{tab:trainer-hooks}).  For~a~dataset bundled with~the~runtime you
implement only \texttt{\_create\_model}: subclass one of~the~ready-made
mixins~--- \texttt{MNISTTrainer}, \texttt{MNISTSequenceTrainer}, or
\texttt{CIFAR10Trainer}~--- which already supply the~other two.

\begin{table}[H]
\centering\small
\begin{tabularx}{\textwidth}{@{}l >{\raggedright\arraybackslash}X@{}}
\toprule
\textbf{Hook (you override)} & \textbf{Signature and~contract} \\
\midrule
\texttt{\_create\_model}     & \texttt{(self) -> nn.Module}.  Return the~model.  \texttt{BaseTrainer} moves it~to~the~device and, on~a~multi-GPU node, wraps it~in~\texttt{nn.DataParallel} automatically. \\
\addlinespace
\texttt{\_load\_datasets}    & \texttt{(self) -> tuple[Dataset, Dataset]}.  Return \texttt{(train, test)} \emph{datasets}~--- not~\texttt{DataLoader}s; \texttt{BaseTrainer} builds the~loaders itself (train: \texttt{BATCH\_SIZE}, shuffled, \texttt{drop\_last}; test: \texttt{min(256, BATCH\_SIZE)}). \\
\addlinespace
\texttt{\_preprocess\_batch} & \texttt{(self, images, labels) -> tuple[Tensor, Tensor]}.  Reshape or~move a~batch before the~forward pass (e.g.\ \texttt{images.squeeze(1)} to~feed an~LSTM); return \texttt{(inputs, labels)}. \\
\bottomrule
\end{tabularx}
\caption{The~three trainer hooks.  With~a~bundled dataset only the~first
is~required~--- the~mixin supplies the~other two.}\label{tab:trainer-hooks}
\end{table}

\paragraph{What you get for~free.}  In~exchange for~honouring the~hooks,
\texttt{BaseTrainer} supplies the~behaviours the~scheduler relies
on~(Table~\ref{tab:job-contract}); the~mechanism behind each is~detailed
in~\S\ref{sec:checkpoints}.

\begin{table}[H]
\centering\small
\begin{tabularx}{\textwidth}{@{}l >{\raggedright\arraybackslash}X >{\raggedright\arraybackslash}X@{}}
\toprule
\textbf{The~image must\ldots} & \textbf{How \texttt{BaseTrainer} does it} & \textbf{Why IJM needs it} \\
\midrule
checkpoint after every epoch & writes \texttt{/checkpoints/latest.pt} via~a~temp-file-and-rename & the~checkpoint is~the~\emph{only} record of~progress, so a~job can~be~killed at~any~instant \\
resume on~start & loads \texttt{/checkpoints/latest.pt} (model, optimiser, epoch) before training & makes (re)dispatch idempotent: the~same image re-run on~another node continues rather than restarts \\
report progress on~stdout & logs \texttt{Epoch x/y - Loss: \ldots\ - Acc: \ldots\% - <s>s} once per~epoch & the~worker parses the~line for~live progress, and~the~trailing seconds is~the~per-epoch compute the~profiler records \\
take config from~the~environment & reads \texttt{EPOCHS\_TOTAL}, \texttt{BATCH\_SIZE} (and~\texttt{CHECKPOINT\_DIR}) & lets the~worker set the~epoch budget per~run~--- \texttt{profilingEpochsNo} for~a~profile, \texttt{epochsTotal} for~a~standard run \\
use all~assigned GPUs & wraps the~model in~\texttt{nn.DataParallel} across the~$N$ GPUs the~worker exposes & a~2-GPU placement only pays off if~the~job actually shards its~batch over both \\
tolerate \texttt{SIGKILL} & no~signal handler~--- a~kill just ends the~process & stopping a~job is~plain \texttt{docker kill}; safety comes from~the~epoch-boundary checkpoint, not~graceful shutdown \\
\bottomrule
\end{tabularx}
\caption{The~training-job contract.  An~image that~ignores it~still~\emph{runs},
but is~not~stoppable or~resumable: a~kill loses all~progress and~the~scheduler
cannot migrate it.}\label{tab:job-contract}
\end{table}

\paragraph{Scope of~the~built-in loop.}  The~supplied loop assumes
\emph{supervised single-label classification}: it~pairs an~\texttt{Adam}
optimiser ($\text{lr}=10^{-3}$) with~\texttt{CrossEntropyLoss}, feeds each
batch as~\texttt{(inputs, labels)}, and~scores top-1 accuracy.  A~workload that
fits this shape needs only the~hooks above.  One that~does~not~--- a~regression
target, a~custom loss, a~language-model objective~--- overrides \texttt{train}
or~\texttt{\_train\_one\_epoch} while still inheriting the~checkpoint and
progress contract, so it~stays stoppable and~resumable.  An~image that~bypasses
\texttt{BaseTrainer} altogether still runs, but~forfeits those guarantees.

\paragraph{A~complete job.}  The~following \texttt{train.py} is~a~full,
runnable IJM job in~its entirety~--- a~CNN on~CIFAR-10, using the~bundled
mixin so only \texttt{\_create\_model} is~written:

\begin{lstlisting}[language=Python, numbers=none]
# train.py
import torch.nn as nn
from base import CIFAR10Trainer        # bundled dataset mixin

class Net(nn.Module):
    def __init__(self, num_classes=10):
        super().__init__()
        self.body = nn.Sequential(
            nn.Conv2d(3, 32, 3, padding=1), nn.ReLU(),
            nn.AdaptiveAvgPool2d(1))
        self.head = nn.Linear(32, num_classes)

    def forward(self, x):
        return self.head(self.body(x).flatten(1))

class Trainer(CIFAR10Trainer):         # data + batch wiring inherited
    def _create_model(self) -> nn.Module:
        return Net()

if __name__ == "__main__":
    Trainer().train()   # stop/resume/progress: automatic
\end{lstlisting}

\paragraph{Packaging and~runtime inputs.}  The~image bundles \texttt{base.py}
beside the~script and~runs it~as~the~entrypoint:

\begin{lstlisting}[language={}, numbers=none]
FROM python:3.13-slim
RUN pip install torch torchvision
COPY base.py train.py ./
CMD ["python", "-u", "train.py"]
\end{lstlisting}

At~dispatch the~worker sets \texttt{EPOCHS\_TOTAL} and~\texttt{BATCH\_SIZE}
(a~profile run passes \texttt{profilingEpochsNo} as~\texttt{EPOCHS\_TOTAL},
Table~\ref{tab:job-create}) and~bind-mounts the~per-job checkpoint directory
at~the~default \texttt{CHECKPOINT\_DIR}\,=\,\texttt{/checkpoints}.
The~legacy P600 node, pinned to~CUDA~10.1 / PyTorch~1.5.1, is~served by~a
parallel \texttt{Dockerfile.legacy} selected through \texttt{IMAGE\_TAG\_OVERRIDE}
(\S\ref{sec:checkpoints}); the~only constraint this~adds is~that~the~trainer
module must \emph{import} under Python~3.7.

\paragraph{Datasets.}  The~bundled mixins fetch MNIST and~CIFAR-10 into~the
shared cache at~\texttt{/runs/data} (Table~\ref{tab:shared-fs}), staged once
and~mounted read-only on~every node.  A~custom dataset must currently be~baked
into~the~image or~added to~that~shared cache: the~\texttt{directoryToMount}
submission field is~retained for~ANDREAS compatibility but~is~not~bind-mounted
by~the~present worker.

\section{Full API overview}

Table~\ref{tab:api-endpoints} lists every endpoint.  Interactive
documentation is~available at~\texttt{/docs} (Swagger~UI) and
\texttt{/redoc} when the~API is~running.

\begin{table}[H]
\centering
\begin{tabularx}{\textwidth}{@{}l l X@{}}
\toprule
\textbf{Method} & \textbf{Path} & \textbf{Description} \\
\midrule
\texttt{GET}    & \texttt{/}                            & Health check endpoint \\
\texttt{GET}    & \texttt{/health}                      & Detailed health check~--- pings DB and~checks job runner \\
\texttt{GET}    & \texttt{/nodes}                       & List all~cluster nodes with~their current status \\
\texttt{GET}    & \texttt{/gpu-costs}                   & Hourly GPU energy costs used by~the~scheduler for~optimisation \\
\texttt{GET}    & \texttt{/configurations}              & List all~valid hardware configurations in~the~cluster \\
\texttt{POST}   & \texttt{/jobs}                        & Create a~new job \\
\texttt{GET}    & \texttt{/jobs}                        & List all~jobs ordered by~creation time (paginated via \texttt{limit}/\texttt{offset}) \\
\texttt{GET}    & \texttt{/jobs/\{id\}}                 & Get a~specific job by~ID \\
\texttt{POST}   & \texttt{/jobs/\{id\}/stop}            & Request job to~be stopped (user-driven; sticky \textsc{preempted} state) \\
\texttt{POST}   & \texttt{/jobs/\{id\}/resume}          & Request job to~be resumed \\
\texttt{DELETE} & \texttt{/jobs/\{id\}}                 & Delete a~job, killing any~running container first \\
\texttt{DELETE} & \texttt{/jobs}                        & Delete all~jobs and~their profiling results \\
\texttt{GET}    & \texttt{/jobs/\{id\}/logs}            & Return the~output log for~a~job by~proxying to~the~worker that~owns it \\
\texttt{GET}    & \texttt{/profiling-results/\{id\}}    & All~profiling results for~a~job, ordered by~duration (fastest first) \\
\texttt{GET}    & \texttt{/admin/slots}                 & NodeSlots snapshot: in-memory \texttt{used} vs.\ DB-authoritative count \\
\texttt{POST}   & \texttt{/admin/reconcile-slots}       & Force NodeSlots to~match DB-authoritative state (idempotent) \\
\texttt{GET}    & \texttt{/admin/dispatch-tasks}        & List in-flight \texttt{\_dispatch\_when\_slot\_free} tasks \\
\bottomrule
\end{tabularx}
\caption{API endpoints.}\label{tab:api-endpoints}
\end{table}

Three architectural choices set IJM's control plane apart from ANDREAS.
First, the~\emph{job-management} logic~--- scheduling, profiling, and~dispatch~---
is~consolidated into a~single FastAPI service, where ANDREAS used a~separate
Java Jobs Manager.  The~cost-aware optimiser itself is~\emph{not}
reimplemented: IJM consumes the~GPUspb
solver~\cite{filippini2023pathrelinking,filippini2024stochastic}~--- the~same
line of~work as~ANDREAS's optimiser~--- as~an~external service, so what~is
consolidated is~the~job manager, not~the~optimiser.
Second, execution lives in~per-node \emph{workers}~--- also FastAPI~---
that the~API reaches over HTTP through SSH tunnels; that same GPUspb
optimiser is~reached the~same way and~owns every standard-run placement
decision.
Third, IJM exposes explicit \texttt{stop}, \texttt{resume},
\texttt{logs}, and~\texttt{delete} endpoints, together with
\texttt{/profiling-results} and~\texttt{/configurations}, so the~user
has direct visibility into the~scheduler's state.

\section{Runtime environment}

Table~\ref{tab:env-vars} lists the~runtime knobs that influence the
scheduler.  Unless noted otherwise, every variable is~read by~the~API
process at~startup.

\begin{table}[H]
\centering
\begin{tabularx}{\textwidth}{@{}l l X@{}}
\toprule
\textbf{Variable} & \textbf{Default} & \textbf{Description} \\
\midrule
\texttt{DATABASE\_URL}                  & \texttt{postgresql://\ldots/ijm} & PostgreSQL DSN, reached via~SSH tunnel \\
\texttt{OPTIMIZER\_URL}                 & unset & The~GPUspb HTTP service used for~standard-run placement \\
\texttt{OPTIMIZER\_METHOD}              & \texttt{RG} & GPUspb solver method passed in~the~request body \\
\texttt{OPTIMIZER\_VERBOSE}             & 0 & Per-call log verbosity (0--3) \\
\texttt{PROFILING\_CONFIGS\_PER\_JOB}   & 2 & How many distinct GPU configurations a~single instance profiles before falling through to~a~standard run \\
\texttt{IJM\_DRIFT\_HEARTBEAT\_S}       & 15 & Period of~the~slot-tracker drift heartbeat that reconciles the~in-memory slot count against the~database-authoritative count per~node \\
\texttt{WORKER\_GPU\_MODE}              & \texttt{runtime} & Worker-side; how Docker exposes GPUs to~training containers (\texttt{runtime} / \texttt{cdi} / \texttt{none}) \\
\texttt{IMAGE\_TAG\_OVERRIDE}           & unset & Worker-side; rewrites \texttt{:latest} to~the~given tag.  Used on~\texttt{matemagician} to~pick the~CUDA-10.1-compatible \texttt{:legacy} image \\
\texttt{HOST\_ROOT}, \texttt{HOST\_PROJECT\_ROOT} & \texttt{/host} & Resolve container-internal vs.~host paths for~Docker bind mounts; \texttt{HOST\_PROJECT\_ROOT} falls back to \texttt{HOST\_ROOT} when unset \\
\bottomrule
\end{tabularx}
\caption{Key environment variables.}\label{tab:env-vars}
\end{table}
